\topic{1}{Lagrangian i oscylator harmoniczny}{
Lagrangian to różnica energii kinetycznej i potencjalnej
$L(q,\dot q,t)=T-U$. Zasada stacjonarnego działania
$S=\int L\,dt$ prowadzi do równania Eulera--Lagrange'a
$d(\partial L/\partial\dot q)/dt-\partial L/\partial q=0$.
Dla masy na sprężynie
$L=m\dot x^2/2-kx^2/2$, zatem
$m\ddot x+kx=0$.

Rozwiązanie ogólne to
$x(t)=A\cos(\omega_0t)+B\sin(\omega_0t)$,
gdzie $\omega_0=\sqrt{k/m}$.
Warunki $x(0)=x_0$ i $\dot x(0)=v_0$
dają $A=x_0$, $B=v_0/\omega_0$.
Wprowadzenie współrzędnych uogólnionych
pozwala uwzględnić więzy bez jawnego
rozpisywania sił reakcji. Ważne jest,
aby w równaniu ruchu uwzględnić siły
niezachowawcze osobno jako siły uogólnione,
jeśli nie można ich zawrzeć w potencjale.}
\topic{2}{Stopnie swobody i ruch na sferze}{
Swobodny punkt w przestrzeni ma trzy stopnie
swobody. Więz $x^2+y^2+z^2=R^2$ usuwa
jeden, pozostawiając dwa: kąty $\theta,\phi$.
Dla punktu o masie $m$ poruszającego się
bez tarcia po nieruchomej sferze
$L=mR^2(\dot\theta^2+\sin^2\theta\dot\phi^2)/2$.

Równanie dla $\phi$ daje
$d(mR^2\sin^2\theta\dot\phi)/dt=0$,
ponieważ $\phi$ jest współrzędną cykliczną.
Dla $\theta$ otrzymujemy
$\ddot\theta-\sin\theta\cos\theta\,\dot\phi^2=0$.
Zachowany pęd sprzężony z $\phi$
jest składową momentu pędu względem
osi $z$. W tym przykładzie siła
reakcji więzu nie wykonuje pracy
w dozwolonych przesunięciach
wirtualnych, dlatego nie pojawia
się w równaniach dla kątów.}
\newpage
\topic{3}{Hamiltonian i oscylator harmoniczny}{
Pęd kanoniczny definiujemy jako
$p=\partial L/\partial\dot q$, a Hamiltonian
$H=p\dot q-L$ po wyeliminowaniu prędkości.
Dla oscylatora $p=m\dot x$ oraz
$H=p^2/(2m)+kx^2/2$.
Równania Hamiltona to
$\dot x=\partial H/\partial p=p/m$
i $\dot p=-\partial H/\partial x=-kx$.

Ich połączenie daje
$\ddot x+\omega_0^2x=0$.
W przestrzeni fazowej trajektorie
o stałej energii są elipsami:
$p^2/(2m)+kx^2/2=E$.
Dla Hamiltonianu niezależnego
jawnie od czasu energia jest
stała. Formalizm Hamiltona
opisuje ewolucję punktu
$(x,p)$, co ułatwia analizę
symetrii, zachowanych wielkości
i przejście do mechaniki
statystycznej oraz kwantowej.}
\topic{4}{Wahadło matematyczne w opisie Hamiltona}{
Jedyną współrzędną wahadła
o długości $\ell$ jest kąt $\theta$.
Dla masy punktowej $m$ i potencjału
liczonego od dolnego położenia
$L=m\ell^2\dot\theta^2/2
-mg\ell(1-\cos\theta)$.
Pęd kanoniczny wynosi
$p_\theta=m\ell^2\dot\theta$.

Hamiltonian jest sumą
$H=p_\theta^2/(2m\ell^2)
+mg\ell(1-\cos\theta)$.
Równania Hamiltona:
$\dot\theta=p_\theta/(m\ell^2)$
i $\dot p_\theta=-mg\ell\sin\theta$.
W przybliżeniu $|\theta|\ll1$ mamy
$\sin\theta\simeq\theta$,
więc $\ddot\theta+(g/\ell)\theta=0$
i okres $T\simeq2\pi\sqrt{\ell/g}$.
Przy dużych amplitudach okres
zależy od amplitudy; przybliżenie
małych drgań przestaje działać.}
\newpage
\topic{5}{Ładunek próbny i siła Lorentza}{
Ładunek próbny jest na tyle mały,
że nie zmienia istotnie badanego pola.
Na ładunek $q$ poruszający się
z prędkością $\mathbf v$ działa
siła Lorentza
$\mathbf F=q(\mathbf E+\mathbf v\times\mathbf B)$.
Stałe jednorodne pole elektryczne
przyspiesza cząstkę w kierunku
$q\mathbf E$; bez pola magnetycznego
tor jest paraboliczny przy
odpowiednich warunkach początkowych.

Samo pole magnetyczne nie wykonuje
pracy, bo $(\mathbf v\times\mathbf B)
\cdot\mathbf v=0$.
Składowa prędkości prostopadła
do $\mathbf B$ daje ruch kołowy
o promieniu $r=m v_\perp/(|q|B)$
w granicy nierelatywistycznej,
a równoległa ruch jednostajny.
Złożenie tych ruchów tworzy
helisę. Znak ładunku odwraca
kierunek zakrzywienia toru.}
\topic{6}{Prawo Gaussa dla pól E i B}{
Prawo Gaussa dla elektryczności
w postaci całkowej to
$\oint_S\mathbf E\cdot d\mathbf S
=Q_{\rm wew}/\varepsilon_0$,
a lokalnie
$\nabla\cdot\mathbf E=\rho/\varepsilon_0$.
Dla pola magnetycznego
$\oint_S\mathbf B\cdot d\mathbf S=0$
i $\nabla\cdot\mathbf B=0$;
strumień wchodzący i wychodzący
z zamkniętej powierzchni
bilansuje się.

Dla punktowego ładunku w próżni
symetria sferyczna daje
$E(4\pi r^2)=q/\varepsilon_0$,
stąd prawo Coulomba
$\mathbf E=q\hat{\mathbf r}/
(4\pi\varepsilon_0r^2)$.
Zastosowanie wymaga, by na
wybranej sferze wartość $E$
była stała i pole radialne.
Prawo magnetyczne opisuje
brak obserwowanych izolowanych
monopoli w klasycznych
równaniach Maxwella.}
\newpage
\topic{7}{Prawo Faradaya i indukowana SEM}{
Prawo Faradaya w postaci całkowej
mówi $\oint_C\mathbf E\cdot
d\boldsymbol\ell=-d\Phi_B/dt$
dla nieruchomego konturu $C$,
gdzie $\Phi_B=\int_S\mathbf B
\cdot d\mathbf S$.
Lokalnie jest to
$\nabla\times\mathbf E
=-\partial\mathbf B/\partial t$.
Zmienne pole magnetyczne
wytwarza wirowe pole elektryczne.

Dla poruszającego się obwodu
SEM zawiera także człon
z siłą magnetyczną
$\mathbf v\times\mathbf B$.
W generatorze obrót ramki
zmienia strumień i wywołuje
okresową SEM. Znak w prawie
Faradaya wyraża regułę Lenza:
prąd indukowany przeciwdziała
zmianie strumienia, która go
wywołała. Rozróżnienie
nieruchomej i ruchomej
granicy całkowania zapobiega
błędnej interpretacji
samego członu $\partial B/\partial t$.}
\topic{8}{Prawo Ampère'a--Maxwella i kondensator}{
Pełne prawo ma postać
$\oint_C\mathbf B\cdot
d\boldsymbol\ell
=\mu_0(I_{\rm przew}
+\varepsilon_0\,d\Phi_E/dt)$.
W postaci lokalnej:
$\nabla\times\mathbf B
=\mu_0\mathbf J
+\mu_0\varepsilon_0
\partial\mathbf E/\partial t$.
Człon z pochodną pola
elektrycznego nazywa się
prądem przesunięcia.

Podczas ładowania kondensatora
prądem stałym przewodzenia $I$
przez przewód, między okładkami
nie płynie taki sam prąd
nośników przez próżnię.
Pole $E$ jednak narasta
i dla idealnego kondensatora
$\varepsilon_0\,d\Phi_E/dt=I$.
Dzięki temu wynik prawa
Ampère'a nie zależy od
wyboru powierzchni rozpiętej
na tym samym konturze.
Poprawka jest zgodna
z lokalnym zachowaniem ładunku.}
\newpage
